\section{Problem formulation}
\label{sec:problem}

\subsection{The monitored quantities and their null hypotheses}
\label{sec:problem-quantities}

\paragraph{VaR exceedances.}
A model forecasts, before day $t$, a loss quantile $\VaR_t$ that the loss $L_t$ of day $t$ should exceed with probability $\alpha$ (a decimal, here $0.01$). The hit of day $t$ is the flag
$I_t=\indic\{L_t>\VaR_t\}$. Let $\mathcal F_t$ be the information available after day $t$. The null hypothesis of correct conditional coverage states that, given everything that happened before the
day, the probability of a hit is the nominal one,
\begin{equation}
H_0^{\mathrm{VaR}}:\quad \Prob\bigl(I_t=1\mid\mathcal F_{t-1}\bigr)=\alpha\quad\text{for every day }t,
\label{eq:null-var}
\end{equation}
which makes the hits independent Bernoulli($\alpha$) variables. The departures a validator looks for are of two kinds: a hit rate $p=m\alpha$ with a multiplier $m>1$ (the VaR understates the risk,
the case a supervisor penalises) or $m<1$ (the VaR is too conservative), and hits that cluster, where the probability of a hit after a hit is $q>\alpha$ while the hit rate stays $\alpha$.

\paragraph{Defaults against the predicted PD of each grade.}
A rating system has $G$ grades. At step $t$ (a year, a vintage age or a quarter) grade $g$ holds $n_{t,g}$ obligors, of which $d_{t,g}$ default, and the model predicts the default probability $p_{t,g}\in(0,1)$. The
null hypothesis of correct calibration states that, given the past, the defaults of a step are binomial with the predicted probability,
\begin{equation}
H_0^{\mathrm{PD}}:\quad d_{t,g}\mid\mathcal F_{t-1}\;\sim\;\mathrm{Binomial}\bigl(n_{t,g},\,p_{t,g}\bigr)\quad\text{for every step }t\text{ and grade }g.
\label{eq:null-pd}
\end{equation}
The departure of interest is a multiplicative miscalibration: the true default probability of grade $g$ is $m\,p_{t,g}$ with $m>1$ (the model is optimistic) or $m<1$. The null has no asset correlation: it
states that the defaults of a step are conditionally independent, which is the assumption of the binomial test of one grade and of the chi-square test across grades. \Cref{sec:theory} shows what happens to
the monitor when that assumption fails.

\subsection{The monitoring problem}
\label{sec:problem-monitoring}

A monitor reads the data as they arrive and, after each step, either continues or rejects the null. Its decision is a stopping time $\tau$ with respect to the filtration $(\mathcal F_t)$, the first step at which the monitor rejects, and $\tau=\infty$ when it never does.
Three properties define a good monitor. The first is the anytime-valid size: the probability under the null that the monitor ever rejects is at most the stated level $a$, however often it is read and
whenever the analyst stops,
\begin{equation}
\Prob_0(\tau<\infty)\le a .
\label{eq:anytime-size}
\end{equation}
The second is the detection delay: under an alternative with a planted departure from step $0$, the distribution of $\tau$, summarised by its median, since a monitor that misses the departure in
more than half of the streams has an infinite median. The third is the power at a fixed horizon $n$, $\pi(n)=\Prob_1(\tau\le n)$, which is what a fixed-window test is judged on, and the power cost of a monitor is the difference between the power of a fixed-window test read once at $n$ and $\pi(n)$ of the monitor.
A rule that is read every day also has a size over a finite horizon $T$, $\Prob_0(\tau\le T)$; for a fixed-window test read repeatedly, that size is what the study measures in
place of the nominal level of one reading.

\subsection{Notation and assumptions}
\label{sec:problem-notation}

\Cref{tab:notation} lists every symbol used in the paper, with its meaning and its unit, so that a reader can follow the equations of the later sections without searching for a definition. Three assumptions bound the results. The first is that the null of \cref{eq:null-var,eq:null-pd} is the correct description of a
well-calibrated model: for VaR it means conditional coverage, for PD binomial defaults with no cross-sectional or serial dependence beyond what the predicted probability absorbs. The second is that the
mixture weights and the alternatives of every monitor are fixed before any outcome is seen, which keeps each process an e-process (\cref{prop:validity}). The third is that the series are observed in time order and
that a missing observation is an error, not a non-event: the monitors refuse a missing value, because treating a missing day as a day without a hit would bias the evidence against the alternative.

\begin{table}[htbp]
\centering
\small
\caption{Notation: each symbol of the paper with its meaning and its unit; all probabilities are decimals, all times are counts of days, steps or years, and the statistics $E_t$ and $\log E_t$ are dimensionless.}
\label{tab:notation}
\begin{tabularx}{\textwidth}{>{\raggedright\arraybackslash}p{2.6cm}>{\raggedright\arraybackslash}X>{\raggedright\arraybackslash}p{3.3cm}}
\toprule
Symbol & Meaning & Unit \\
\midrule
$t,\;T,\;n$ & Step index, horizon of a run, horizon of one fixed-window reading (250 days) & Count \\
$\alpha$ & VaR exceedance probability, the level of the VaR (0.01) & Probability \\
$I_t$ & Hit flag of day $t$, $\indic\{L_t>\VaR_t\}$ & 0 or 1 \\
$\mathcal F_t$ & Information available after step $t$ & Sigma-field \\
$a$ & Size of a monitor (0.05); it rejects at $E_t\ge 1/a=20$ & Probability \\
$\tau$ & First step at which a monitor rejects ($\infty$ if never) & Count \\
$m,\;\mathcal M,\;K$ & Multiplier of the null parameter, the grid of multipliers of a monitor, its size & Dimensionless, count \\
$w_k$ & Mixture weight of the $k$-th alternative, fixed before the data & Dimensionless \\
$E_t,\;\log E_t$ & E-process of a monitor after step $t$ and its natural logarithm & Dimensionless \\
$\Lambda_{t}(p)$ & Likelihood ratio of a Bernoulli($p$) hit rate to the null after $t$ days & Dimensionless \\
$q$ & Probability of a hit the day after a hit & Probability \\
$G,\;g$ & Number of rating grades, grade index & Count \\
$n_{t,g},\;d_{t,g}$ & Obligors and defaults of grade $g$ at step $t$ & Count \\
$p_{t,g},\;\bar p_g$ & Predicted PD of grade $g$ at step $t$, its mean over the steps & Probability \\
$\rho$ & Asset correlation of the one-factor default model & Dimensionless \\
$\mathrm{KL}(p\Vert r)$ & Kullback-Leibler divergence of Bernoulli($p$) from Bernoulli($r$) & Nats per observation \\
$\pi_R(n)$ & Power by day $n$ of rule $R$, the share of streams with a departure it has rejected by $n$ & Probability \\
\bottomrule
\end{tabularx}
\end{table}
